\section{Finite M\"obius slopes and prime-product lifts}
\label{sec:tpc33-mobius-corollaries}

We derive two consequences of the exact second moment.  The first is an
unconditional almost-all-pair theorem for a finite M\"obius model.  The
second verifies that a complete product slope \(a=\ell j\) does not lose
projective coverage merely because \(\ell\) is restricted.

\subsection{A finite M\"obius model}

For a prime \(q\), define \(\mu_q:\F_q\to\R\) by
\begin{equation}
 \mu_q(0)=0,\qquad
 \mu_q(n)=\mu(n),\qquad 1\le n<q,
 \label{eq:tpc33-finite-mobius-model}
\end{equation}
using the standard representatives.  Let
\begin{equation}
 M(x)=\sum_{n\le x}\mu(n).
 \label{eq:tpc33-Mertens-function}
\end{equation}
Then
\begin{equation}
 \overline{\mu_q}=\frac{M(q-1)}q,
 \qquad
 \|\mu_q^\circ\|_2^2\le q.
 \label{eq:tpc33-mobius-mean-energy}
\end{equation}

\begin{theorem}[Almost every finite determinant-fiber M\"obius pair]
\label{thm:tpc33-mobius-almost-all-slopes}
Fix \(0<\eta<1/2\).  For every prime \(q\) and every
\(h\in\F_q^\times\), outside a set
\(\cE_{q,h,\eta}\subset(\F_q^\times)^2\) satisfying
\begin{equation}
 |\cE_{q,h,\eta}|\le q^{2-2\eta},
 \label{eq:tpc33-mobius-exception-set}
\end{equation}
one has
\begin{equation}
 \boxed{
 \cC_{a,s}^{(h)}(\mu_q,\mu_q)
 =\frac{M(q-1)^2}{q}+O(q^{1/2+\eta}).}
 \label{eq:tpc33-mobius-almost-all-formula}
\end{equation}
Consequently,
\begin{equation}
 \cC_{a,s}^{(h)}(\mu_q,\mu_q)=o(q)
 \label{eq:tpc33-mobius-almost-all-cancellation}
\end{equation}
for a proportion \(1-o(1)\) of coefficient pairs \((a,s)\) as
\(q\to\infty\).
\end{theorem}

\begin{proof}
Apply \cref{cor:tpc33-exceptional-projective-directions} with
\(f=g=\mu_q\) and \(\Lambda=q^{1/2+\eta}\).  By
\eqref{eq:tpc33-mobius-mean-energy}, the total variance is at most \(q^3\),
so the exceptional count is at most
\(q^3/q^{1+2\eta}=q^{2-2\eta}\).  The centered main term is
\[
 q\overline{\mu_q}^{\,2}=\frac{M(q-1)^2}{q},
\]
which proves \eqref{eq:tpc33-mobius-almost-all-formula}.  The prime number
theorem is equivalent to \(M(x)=o(x)\), so this main term is \(o(q)\); the
error is also \(o(q)\) because \(\eta<1/2\).
\end{proof}

\begin{remark}[Model boundary]
The map \(n\mapsto\mu_q(n)\) is not a periodic multiplicative function on
the integers; it is the finite vector obtained by truncating the ordinary
M\"obius function before \(q\).  The theorem is therefore a modular model
for determinant-fiber coefficient averaging, not a theorem about one
integer affine fiber.
\end{remark}

\subsection{Complete and incomplete product slopes}

Let \(\cL\subset\F_q^\times\) be arbitrary.  For a fixed
\(\ell\in\cL\), multiplication by \(\ell\) permutes
\(\F_q^\times\).

\begin{proposition}[Exact product-slope lift]
\label{prop:tpc33-product-slope-lift}
For every \(f,g:\F_q\to\C\),
\begin{align}
 &\boxed{
 \sum_{\ell\in\cL}
 \sum_{j,s\in\F_q^\times}
 \left|
  \cC_{\ell j,s}^{(h)}(f,g)-q\bar f\bar g
 \right|^2}
 \notag\\[-0.2em]
 &\boxed{\qquad
 =|\cL|
 \sum_{a,s\in\F_q^\times}
 \left|
  \cC_{a,s}^{(h)}(f,g)-q\bar f\bar g
 \right|^2.}
 \label{eq:tpc33-product-slope-lift}
\end{align}
In particular, restricting \(\cL\) to prime residue classes causes no
coverage loss when \(j\) runs through a complete nonzero residue system.
\end{proposition}

\begin{proof}
For each fixed \(\ell\ne0\), the substitution \(a=\ell j\) is a
bijection of \(\F_q^\times\).  Apply it inside the \(j\)-sum and then sum
over \(\ell\).
\end{proof}

For an incomplete integer set \(\cJ\), let
\begin{equation}
 \nu_{\cJ}(r)=\#\{j\in\cJ:j\equiv r\pmod q\},
 \qquad
 \nu_{\cJ}^*=\max_{r\in\F_q^\times}\nu_{\cJ}(r),
 \label{eq:tpc33-residue-multiplicity}
\end{equation}
and define \(\nu_{\cS}^*\) similarly.  We restrict to elements nonzero
modulo \(q\).

\begin{corollary}[Incomplete product-slope energy]
\label{cor:tpc33-incomplete-product-slope}
For finite integer sets \(\cJ,\cS\) with no multiples of \(q\),
\begin{align}
 &\sum_{\ell\in\cL}\sum_{j\in\cJ}\sum_{s\in\cS}
 \left|\cC_{\ell j,s}^{(h)}(f,g)-q\bar f\bar g\right|^2
 \notag\\
 &\qquad\le
 |\cL|\nu_{\cJ}^*\nu_{\cS}^*
 q\|f^\circ\|_2^2\|g^\circ\|_2^2.
 \label{eq:tpc33-incomplete-product-slope-energy}
\end{align}
If \(\cJ,\cS\) are intervals of respective lengths \(J,S\), then
\begin{equation}
 \nu_{\cJ}^*\le J/q+1,
 \qquad
 \nu_{\cS}^*\le S/q+1.
 \label{eq:tpc33-interval-residue-multiplicity}
\end{equation}
After division by \(|\cL|JS\), the rms loss relative to complete residue
averaging is at most
\begin{equation}
 \boxed{
 \left(1+\frac qJ\right)^{1/2}
 \left(1+\frac qS\right)^{1/2}.}
 \label{eq:tpc33-incomplete-rms-loss}
\end{equation}
\end{corollary}

\begin{proof}
Group \(j,s\) by their nonzero residues.  Each pair is repeated at most
\(\nu_{\cJ}^*\nu_{\cS}^*\) times.  Apply
\cref{prop:tpc33-product-slope-lift} and then
\eqref{eq:tpc33-projective-upper-bound}.  The interval bounds are
elementary.  Dividing the resulting mean square by the complete-residue
mean-square normalization introduces an additional factor
\((q-1)/q\le1\); discarding it gives
\eqref{eq:tpc33-incomplete-rms-loss}.
\end{proof}

The exact lift is useful only after the orbit variable covers residue
classes economically.  When \(q\gg J\), the first factor in
\eqref{eq:tpc33-incomplete-rms-loss} is large; this is the coverage side of
the single-modulus obstruction proved later.
